%%
%% This is file `sample-sigconf.tex',
%% generated with the docstrip utility.
%%
%% The original source files were:
%%
%% samples.dtx  (with options: `all,proceedings,bibtex,sigconf')
%% 
%% IMPORTANT NOTICE:
%% 
%% For the copyright see the source file.
%% 
%% Any modified versions of this file must be renamed
%% with new filenames distinct from sample-sigconf.tex.
%% 
%% For distribution of the original source see the terms
%% for copying and modification in the file samples.dtx.
%% 
%% This generated file may be distributed as long as the
%% original source files, as listed above, are part of the
%% same distribution. (The sources need not necessarily be
%% in the same archive or directory.)
%%
%%
%% Commands for TeXCount
%TC:macro \cite [option:text,text]
%TC:macro \citep [option:text,text]
%TC:macro \citet [option:text,text]
%TC:envir table 0 1
%TC:envir table* 0 1
%TC:envir tabular [ignore] word
%TC:envir displaymath 0 word
%TC:envir math 0 word
%TC:envir comment 0 0
%%
%%
%% The first command in your LaTeX source must be the \documentclass
%% command.
%%
%% For submission and review of your manuscript please change the
%% command to \documentclass[manuscript, screen, review]{acmart}.
%%
%% When submitting camera ready or to TAPS, please change the command
%% to \documentclass[sigconf]{acmart} or whichever template is required
%% for your publication.
%%
%%
\documentclass[sigconf]{acmart}

%% Rights management information.  This information is sent to you
%% when you complete the rights form.  These commands have SAMPLE
%% values in them; it is your responsibility as an author to replace
%% the commands and values with those provided to you when you
%% complete the rights form.
\setcopyright{acmlicensed}
\copyrightyear{2018}
\acmYear{2018}
\acmDOI{XXXXXXX.XXXXXXX}

%% These commands are for a PROCEEDINGS abstract or paper.
\acmConference[ICSE '25]{IEEE/ACM International Conference on Software Engineering}{April 27--May 3,
  2025}{Ottawa, Ontario, Canada}
%%
%%  Uncomment \acmBooktitle if the title of the proceedings is different
%%  from ``Proceedings of ...''!
%%
%%\acmBooktitle{Woodstock '18: ACM Symposium on Neural Gaze Detection,
%%  June 03--05, 2018, Woodstock, NY}
\acmISBN{978-1-4503-XXXX-X/18/06}


%%
%% Submission ID.
%% Use this when submitting an article to a sponsored event. You'll
%% receive a unique submission ID from the organizers
%% of the event, and this ID should be used as the parameter to this command.
%%\acmSubmissionID{123-A56-BU3}

%%
%% For managing citations, it is recommended to use bibliography
%% files in BibTeX format.
%%
%% You can then either use BibTeX with the ACM-Reference-Format style,
%% or BibLaTeX with the acmnumeric or acmauthoryear sytles, that include
%% support for advanced citation of software artefact from the
%% biblatex-software package, also separately available on CTAN.
%%
%% Look at the sample-*-biblatex.tex files for templates showcasing
%% the biblatex styles.
%%

%%
%% The majority of ACM publications use numbered citations and
%% references.  The command \citestyle{authoryear} switches to the
%% "author year" style.
%%
%% If you are preparing content for an event
%% sponsored by ACM SIGGRAPH, you must use the "author year" style of
%% citations and references.
%% Uncommenting
%% the next command will enable that style.
%%\citestyle{acmauthoryear}


%%
%% end of the preamble, start of the body of the document source.
\begin{document}

%%
%% The "title" command has an optional parameter,
%% allowing the author to define a "short title" to be used in page headers.
\title{Visualising Developer Interactions in Code Reviews}

%%
%% The "author" command and its associated commands are used to define
%% the authors and their affiliations.
%% Of note is the shared affiliation of the first two authors, and the
%% "authornote" and "authornotemark" commands
%% used to denote shared contribution to the research.
\author{Daniel Bee}
\affiliation{%
  \institution{Royal Holloway, University of London}
  \city{Egham}
  \country{UK}
}
\email{Daniel.Bee.2022@live.rhul.ac.uk}

\author{DongGyun Han}
\affiliation{%
  \institution{Royal Holloway, University of London}
  \city{Egham}
  \country{UK}
}
\email{DongGyun.Han@rhul.ac.uk}


%%
%% By default, the full list of authors will be used in the page
%% headers. Often, this list is too long, and will overlap
%% other information printed in the page headers. This command allows
%% the author to define a more concise list
%% of authors' names for this purpose.
\renewcommand{\shortauthors}{Bee et al.}

%%
%% The abstract is a short summary of the work to be presented in the
%% article.
\begin{abstract}
  The act of code review (also known as a pull request) is a code change auditing technique performed by developers that are not authors of the original code change(s).
    Recent studies have demonstrated diverse benefits of code review.
    For example, Bacchelli and Bird reported that code review is effective to share knowledge between developers to improve code changes. % (REFERENCE HERE)
    This paper aims to outline how this operation can be assembled into an asynchronous, seamless visualisation that can be used to gauge the relationships between different developers in a given project.
    The development of this final product was split into three stages:

    \begin{enumerate}
        \item Mining code review comments from open-source production code on GitHub.
        \item Analysing interactions between developers using a migration script to a backend relational MariaDB database.
        \item Visualising the interactions between developers using D3.js.
    \end{enumerate}

    These three stages were amalgamated into a full-stack web service that took advantage of several web technologies that have transcended the backend and frontend development landscapes.
\end{abstract}

%%
%% The code below is generated by the tool at http://dl.acm.org/ccs.cfm.
%% Please copy and paste the code instead of the example below.
%%
\begin{CCSXML}
<ccs2012>
   <concept>
       <concept\_id>10011007.10011074.10011111.10011696</concept\_id>
       <concept\_desc>Software and its engineering~Maintaining software</concept\_desc>
       <concept\_significance>500</concept\_significance>
       </concept>
 </ccs2012>
\end{CCSXML}

\ccsdesc[500]{Software and its engineering~Maintaining software}

%%
%% Keywords. The author(s) should pick words that accurately describe
%% the work being presented. Separate the keywords with commas.
\keywords{code review, visualisation, developer interactions}


%%
%% This command processes the author and affiliation and title
%% information and builds the first part of the formatted document.
\maketitle

\section{Introduction}
\label{sec:introduction}

Visualising complex and congested data is a challenge,
no matter how well-versed a developer is on the knowledge of a certain API(s).
Therefore,
to help aid in this process, several APIs and frameworks can be used to help make it easier and more efficient.
For instance, Spring Boot was used for the backend of the project.
This was chosen
because it makes programming REST APIs incredibly simple with its model that is built on several levels of abstraction.
In addition to this,
Spring Boot is powered by Java,
which is a programming language that is universally applied in by millions of developers all over the world.
As a result, the project is more accessible to contributors.
On the other hand, D3.js, React, and Material UI were used in the frontend of the project, alongside GraphQL\@.
These APIs and frameworks are chosen because of their simplicity and professional look for end-users.
Firstly, GraphQL supplements the crawling process to collect the necessary data from GitHub.
Without this,
the backend relational database that stores necessary pull request information would contain no data and migration would not be possible.
Furthermore, React and Material UI go hand-in-han for a polished and high-quality website design.
Moreover,
these libraries give both contributors and developers an immense array of tools for lots of different styles of projects.
A vital part in this project was a npm (Node Package Manager) library: Axios.
Axios is one of the most popular promise-based HTTP libraries that is used in Node.js and web development in general.
This allowed a callback to be made to the backend service to allow the frontend to transform the JSON-based input from the backend into a function D3.js graph.
Finally,
D3.js was paramount
for meeting the specifications for this project
since it encompassed the visualisation representing the relationships between developers.
D3.js is incredibly efficient to use, and yet so powerful in what it can achieve.
It was able to support thousands of nodes asynchronously whilst reading from a database from the backend in a matter of seconds.
Lastly, testing was a crucial part of the development cycle that was incorporated from the very beginning.
This coding technique is known as TDD (Test-Driven Development) and helped amend issues early on by testing features in isolation to one another.
JUnit Jupiter was the library chosen to do this, alongside some parts of the Spring Boot Testing framework.
To summarise, the following technologies were chosen to complement the design and development of the project:

\begin{enumerate}
    \item GraphQL
    \item React
    \item Material UI
    \item Axios
    \item D3.js
    \item Spring Boot
    \item JUnit Jupiter
\end{enumerate}

\newpage

\section{Background}
\label{sec:background}


\section{Visualisation}
\label{sec:design}

\section{Related Work}
\label{sec:relatedWork}


\section{Conclusion}
\label{sec:conclusion}

%%
%% The next two lines define the bibliography style to be used, and
%% the bibliography file.
\bibliographystyle{ACM-Reference-Format}
\bibliography{reference}

\end{document}
\endinput
%%
%% End of file `sample-sigconf.tex'.
